


Generated by \texttt{python -m crashguard.rd\_assumptions} from the constants in the code. \textbf{Source "assumption" means: chosen by the project, not taken from an opened document or measured.} Only the Semtech time-on-air formula is a published formula
(reproduced in \texttt{lora.py}; the datasheet itself was not re-opened in the redesign round, so it is listed as "formula as coded"). Every paper table that rests on the simulator points here. Nothing in this table was measured on hardware.

{\scriptsize\raggedright
\begin{longtable}{p{0.237\linewidth}p{0.237\linewidth}p{0.237\linewidth}p{0.237\linewidth}}
\toprule
Constant & Value & Source & Results that depend on it \\
\midrule
\endhead
Path-loss exponent & 3.0 & assumption & world\_sim geometry; all delivery results (\texttt{pa.*}, \texttt{ar.r2.*}, \texttt{rd.*}) \\
Path loss at 1 m & 25.0 dB & assumption & same \\
Transmit power & 14.0 dBm & assumption & same \\
Transmit power, reading of the 10 dB extra loss (r7b G0.2) & the author chose the conservative reading: the 10 dB is path margin, so a compliant dipole node transmits at most 10 dBm; the stored results at 14 dBm assume up to 4.00 dB more margin (about 1.4 spreading-factor steps at 2.8 dB per step) & assumption, ANALYTICAL note & link results (who is in range, delivery): optimistic bound on margin; protocol, security and duty-cycle results: none \\
Extra vehicle and body loss & 10.0 dB & assumption & same \\
Log-normal shadowing (static, per vehicle) & 6.0 dB & assumption & same \\
Cell radius, one gateway at the centre & 3000 m & assumption & same \\
Per-transmission fading & 2.0 dB & assumption (capacity.py) & collision outcomes: \texttt{pa.*}, \texttt{rd.*} \\
Per-vehicle level spread (no-geometry mode) & 8.0 dB & assumption (capacity.py) & earlier studies without geometry \\
Capture rule & a frame survives an overlap if its level exceeds each overlapping frame by 6.0 dB; same SF and channel only (spreading factors are treated as orthogonal) & assumption & every collision count: \texttt{pa.*}, \texttt{rd.d0.c.*}, \texttt{rd.*} \\
Gateway sensitivity SF7 / SF8 / SF9 / SF10 / SF11 / SF12 & -123.0 / -126.0 / -129.0 / -132.0 / -134.5 / -137.0 dBm & assumption (typical 125 kHz values, not from an opened document) & who is in range: all geometry results \\
Loss reduction per SF step above the base SF & x 0.7 on the per-direction loss (uplink and downlink) & assumption (N16); the D1 table reuses it & SF-raise policies; D1 and D2 (\texttt{rd.*}); the simulator-dependence of D1 is tested in the cross-simulator check \\
Per-direction loss levels p & design 0.1, 0.3; held-out 0.2, 0.4; DEV 0.15, 0.35, 0.55; FRESH-A drawn after the freeze around 0.15, 0.30, 0.45, 0.60 & stated stress values & all delivery results \\
Duty limit & 1\% and 0.1\% (old grids); DEV 1\% and 0.2\%; FRESH-A 0.5\% and 10\%; simulator default 1\% (\texttt{s11\_sct.DEFAULT}) & assumption; the Indian rules for 433.05 to 434.79 MHz (G.S.R. 347(E), 2022, rule 3 row 1) allow 10 mW e.r.p. at a duty cycle of at most 10\% with the observation interval not stated, so 1\% is stricter (\texttt{docs/03\_results/FINALIZE\_INDIA\_RULES.md}) & airtime, duty exceed checks \\
Time on air & Semtech SX127x formula, 125 kHz, coding rate 4/5, 8-symbol preamble, explicit header, CRC & formula as coded in \texttt{lora.py} & all airtime results \\
Alert / ACK / flash / heartbeat length & 97 B / 26 B / 14 B / 16 B & project packet format (\texttt{docs/01\_protocol/PACKET\_FORMAT.md}) & airtime; R2 varies the alert length \\
Cluster ACK length (D2) & 11 B header + 2 B per id + 64 B Ed25519 signature; up to 24 ids & design & D2 airtime and deafness \\
ACK processing delay / listen margin & 0.1 s / 0.1 s & assumption & retry timing \\
Compute time before the initial send & 0.05 s & assumption (not measured on the ESP32) & latency results \\
Cancel window by class (Low, Moderate, Severe, Critical) & 10 / 8 / 6 / 5 s & assumption (\texttt{docs/00\_overview/ARCHITECTURE.md} range 5 to 10 s) & severity-table delay before the initial send \\
Class deadlines (Severe, Critical) & 30 s / 10 s from the crash instant & assumption (V7 convention, reused) & all reach-within-deadline metrics; D1 table \\
Burst & b vehicles (Critical) crash within a 10 s window starting at 600 s; b = 20 (design), 30 (held-out), 25 (DEV), drawn 18 to 26 (FRESH-A) & assumption & burst results \\
Crash alert rate & 0.1 per vehicle per hour & stress value, not a real crash rate & steady-state results \\
Event class pool & real out-of-fold class pairs of the VZCrash cross-validation (no synthetic pulse) & MEASURED\_SOFTWARE pool & class assignment, mis-triage \\
Decision time pool & early-exit decision times of the de-duplicated test events (seed 42, 5\% target) & MEASURED\_SOFTWARE pool & latency \\
Simulated hours per cell & old grids 30; DEV 100; FRESH-A 600 & design & interval widths \\
Overhear model & a vehicle decodes a gateway ACK with probability 1 - p x SF decay and when its level at the gateway is above sensitivity & assumption & load index of D1, gate of D2, OS-DR \\
D1 prior & failure mean 0.3, strength 6, forgetting 0.97 & design (tunable on DEV) & D1 \\
D1 retry gaps / SF steps / loss bins / load bins & (0.0, 0.25, 0.5, 1.0, 2.0, 3.5) s / (0, 1, 2) / 8 / 4 & design & D1 \\
D2 defer interval (Severe, Critical) & 4 s / 1 s, at most 3 per alert & design (scale tunable on DEV) & D2 \\
D3 slots (session 2) & slot = alert time on air + 0.02 s guard; epoch at most 5 s (but at least 2 slots); S = clip(ceil(N\_hat / m), 2, 12); thr3 2, m 1, alternate SF on (DEV iteration 3); clock offset at most 5 ms; slot index = SHA-256(epoch, vehicle id) mod S & design; thr3, m, alt tuned on DEV; the others are design constants (S\_MAX = 12 written before any DEV run) & D3 (\texttt{rd.fb.*}); the epoch cap does not hold at SF12 (V12 10.6) \\
Vehicle-to-vehicle overhearing for D3 (session 2) & a vehicle with a pending alert decodes another vehicle's alert or flash if the node-to-node link (path loss and shadowing as for the gateway) is above sensitivity and it is not transmitting; no collisions at vehicle receivers & assumption & D3 burst-size estimate (rarely reached three vehicles in a DEV burst hour) \\
Two-stage frames (D4) & provisional flash 14 B (R3 as registered), 24 B with an 8 B tag (D4), 97 B (task wording); confirm 97 B; CANCEL 14 B x 3; at most 4 tries; exchange success (1 - p)\textasciicircum{}2 & design (R3 values); the R3 events are the stored test events at the 1\% target & H-D4a (\texttt{rd.fb.D4a.*}) \\
ns-3 scenario S-X (session 2) & gateway at 15 m, vehicles at 1.2 m, same channel constants as the custom simulator, sensitivities patched to the custom values, RX1 0.1 s and RX2 0.2 s (policies p1, quick, d1) or 1 s and 2 s (standard LoRaWAN), ACK frame about 13 B, no flash frame, one alert per vehicle & assumption; local ns-3 changes in \texttt{XSIM\_BUILD.md} & cross-simulator verdicts (\texttt{rd.x13.*}) \\
Firmware path (D6) & host build, double precision (bit-exact with Python); ESP32 uses a single-precision FPU, so a port needs float or fixed point (NOT\_RUN) & scope & \texttt{rd.d6.*}; no timing claim \\
Nessa-style baseline (ADDED-AFTER-SESSION-1) & ACK timeout Ub(0) = k x time on air of the 97-byte alert (k = 1 after DEV), Lb(0) = Ub(0) / 2, doubling per stage; redundant copy after the initial retransmission of each stage with spacing U(0, A), A = 1 s after DEV; SF + 1 per stage; 12 tries; no classes, no flash & assumption: from a secondary description (Chalmers thesis, sections 3.1 and 5.3) of Nessa et al. CNS 2020, which was not opened; k and A tuned on DEV; Lb(0) is not in the description & \texttt{rd.fc.*}, \texttt{rd.x13.*} (Nessa rows) \\
Calibration sensitivity settings (A3) & (a) -123, -126, -129, -132, -134.5, -137 dBm; (b) -130, -132.5, -135, -137.5, -140, -142.5 dBm (ns-3 shipped gateway values; ns-3 end-device table set equal in (b)); ADR backoff off in both & both fixed in V12 11.3 before the runs & \texttt{rd.x1.*}, \texttt{rd.x1.b.*} \\
Hostile conditions & none in the delivery simulator (no jammer, no spoofing); security is checked in the bounded models & scope & H-D2c \\
\bottomrule
\end{longtable}}

\texttt{LO.DUTY} (the 10\% allowance of \texttt{lora.py}) is 10\%, which equals the limit of rule 3 row 1 of G.S.R. 347(E) of 2022 (full text read on 2026-11-16; the observation interval is not stated). Regulatory duty limits are parameters, never claims.


